\documentclass[11pt, draft]{article}
\topmargin=0mm
\oddsidemargin=0mm
\textwidth170mm
\textheight230mm
\begin{document}
%\pagestyle{empty}
\language=0
%\hyphenation{}

\begin{center}
{\bf
IMPLEMENTATION OF THE DIFFERENTIAL INVOLUTIVE ALGORITHMS IN THE CAS MAPLE VR5}\\
\vspace {0.3cm}
{Vladimir A.Mityunin}\\
{\it \small Department of Mathematics and Mechanics, Moscow State University, 119899 Moscow,
Russia
\\ e-mail: vmit@park.ru}
\end{center}
\vspace {1.0cm}
Involutive bases are Groebner bases, though, generally redundant,
and involutive methods provide an alternative approach to
computation of Groebner bases. Using involutive basis, one can
easily obtain, for example, dimensional polynomial of the system.
Involutive bases may be used in the Lie symmetries analysis of the
differential systems. In the paper \cite{Ge} generalization of the
involutive basis definition for differential system was given.
Then, in the recent work of Prof. V. Gerdt \cite{DiffCompletion}
an algorithm for construction of the minimal differential
involutive Groebner basis for system of the linear differential
equations was introduced. The main goal of my work was to
implement this algorithm in the CAS Maple. Some of the
optimizations, described in the \cite{DiffCompletion,
Optimization, Mathematica}, were implemented too. The next step
was to perform comparison with the {\it diffalg} package. My
implementation was used in computing of the primitive element in
the differential module for some systems of linear differential
equations and computing of the differential dimensional
polynomials. Unfortunately, Maple is not very effective system,
and we can't use it for computing of the really interesting
examples. But we can analyze and use results in the
implementation on the more low-level language, such as C++. It
seems, that correct implementation of the involutive algorithms
may be much more faster, then any implementation of the classic
Buchberger algorithm.

\begin{thebibliography}{99}
\bibitem{Ge}
\emph{Vladimir P. Gerdt} "Groebner bases and involutive methods
for algebraic and differential equations" // Mathematics and
Computers in Modelling, 25, No. 8/9, 1997, 75-90

\bibitem{GeBl2}
\emph{Vladimir P. Gerdt, Yuri A.Blinkov} "Minimal involutive
bases" // Mathematics and Computers in Simulation 45, 1998,
543-560

\bibitem{DiffCompletion}
\emph{Vladimir P. Gerdt} "Completion of Linear Differential
System to Involution" // Computer Algebra in Scientific Computing,
V.G.Ganzha, E.W.Mayr and E.V. Vorozhtsov (Eds.), Springer-Verlag,
Berlin, 1999, pp.115-137.

\bibitem{Optimization}
\emph{Vladimir P. Gerdt} "Involutive division technique: some
generalization and optimizations" // Zapiski Nauchnykh Seminarov
POMI (St.Petersburg) 258 (1999) 185-206. To appear in J.Math.Sci.
258 (2000).

\bibitem{Mathematica}
\emph{Vladimir P. Gerdt, Matthias Berth and G\"{u}nter Czichowski}
"Involutive Divisions in Mathematica: Implementation and Some
Applications" // Proceedings of the 6th Rhein Workshop on
Computer Algebra (Sankt-Augustin, Germany, March 31 - April 3,
1998), J.Calmet (Ed.), Institute for Algorithms and Scientific
Computing. GMD-SCAI, Sankt-Augustin, 1998, 74-91 .

\bibitem{Kolchin}
\emph{E.R.Kolchin} "Differential algebra and algebraic groups",
New York, Academic Press, 1973

\bibitem{Dif}
\emph{Kondratieva M.V., Levin A.B., Mikhalev A.V., Pankratiev
E.V.} Differential and Difference Dimension Polynomial.  Kluwer
Acad. Publ.,1999.

\end{thebibliography}
\end{document}
